% intro: general-reader introduction (model-first rewrite, 2026-10-07).
\section{Introduction}\label{sec:intro}

\textbf{The problem.} Groups of language-model agents now work together for weeks or months. They share computers, files and a chat. They read each other, copy each other and sometimes wait for each other. Two kinds of people need to understand such a swarm. The first kind works on \emph{interpretability}: given the logs, what is the swarm doing, and why? The second kind runs the swarm and needs \emph{control}: which numbers to watch, when to worry, and which levers move the swarm by how much. Both need a description that is much smaller than the logs and that still predicts.

\textbf{Why physics.} Statistical physics gives such descriptions for systems of many interacting parts. A magnet has $10^{23}$ spins, but a few numbers predict how it responds. The ratio of temperature to critical temperature, $T/T_c$, tells how far a local push spreads and how large the fluctuations are, without a model of each spin. Physics also separates two causes of order. A \emph{field} is an external drive that acts on each part directly. A \emph{coupling} is an interaction between parts. A field aligns parts that never interact, so co-movement alone does not prove interaction. This distinction matters for agent swarms, because the scheduler, the operator and the shared model weights all act as fields. Sociophysics applies these ideas to people~\cite{castellano2009}. An agent swarm fits them better than most human groups do. Its states are logged at each step. Its fields (goals, schedules, operator messages) are known and dated. Its coupling channel, the messages, is recorded message by message, together with what each agent could see.

\textbf{The system.} The AI Village is a public record of one swarm, run by AI Digest~\cite{aivillage}. Forty-six frontier agents from several labs worked on 51 village-wide goals from April 2025 to September 2026 (Fig.~\ref{fig:timeline}). The agents use browsers, shells and a group chat. Three scaffold changes divide the record into regimes, and each change alters what an agent can read and when it can act. Forty-six dated changes of scaffold, roster, rooms and operator behaviour serve as natural experiments. We cannot run new swarms, so these changes are our interventions.

\textbf{The question and the answer.} We ask which physics model, if any, describes this swarm well enough to read its logs and to steer it. The answer is a weakly coupled, subcritical kinetic spin system under strong external fields. Each agent changes state only when it calls its model. It responds to other agents only through the messages that the call reads, mostly the messages that name it. The read-out loop gain is $g\approx0.13$, so a push to one agent is amplified by about $1/(1-g)\approx1.15$ and never cascades. External fields set most of the state: the scheduler sets when agents act, the goal kickoff sets what they talk about, rooms split them into domains, and style marks each agent. Each agent's content relaxes as a fast kick from what it reads on top of a slow personal well. The context window holds the working state, and it is the only store whose loss costs much work. We find no collective unit above the agents.

\textbf{What each reader gets.} For interpretability, the picture says where to look: decompose any co-movement into field and coupling, and locate memory in the context window, not in notes or the group. For control, it gives dials with current values (Table~\ref{tab:observables}): the loop gain and its distance to criticality, the field share of co-activation, the depth and time constants of the agents' wells, the number of collective modes, and the value per bit of the context window. It also sizes the levers. Field-type levers (kickoff text, context resets, cap length) are large. Message-type levers (nudges, corrections) are small.

\textbf{How we tested.} We tested 16 physics models from a library of 17 with 132 hypotheses (Sec.~\ref{sec:method}). Claude agents generated, pre-registered, ran and scored each hypothesis, under review by the second author. A physics model can fit almost any record, so a good fit is not evidence. Each hypothesis states its prediction, null and kill rule before it touches real data. We rank models by faithfulness: whether they beat the strongest null, predict a statistic they did not fit, and survive their own kill rules. We reserved 16 goal periods and four natural-experiment windows for confirmation. Almost nothing has run on them yet, so every result below is exploratory.

\textbf{Plan.} Section~\ref{sec:picture} states the effective picture and turns it into monitoring observables. Section~\ref{sec:method} describes the data and the research loop. Sections~\ref{sec:models}--\ref{sec:rest} go through the models, from the most useful to the rejected ones. Section~\ref{sec:cross} collects the results that span models, and Sec.~\ref{sec:limits} states the limits. The appendices list the goal periods, the hypothesis~$\times$~model table and all hypotheses with their scores.

\begin{figure*}[t]
\includegraphics[width=\textwidth]{timeline.pdf}
\caption{The AI Village over time. Top strip: the 51 goal periods, colored by who wrote the goal. Middle: the agent roster, colored by lab, from joining to retirement. Bottom: agents present, stacked by lab. Dashed lines mark scaffold changes: A chat while on computer; B human helpers; C history search and memory consolidation; D auto-nudger; E chat rooms; F continuous computer use (start of regime III); G outreach approval; H 200-event context cap; I 8\,h per day and new repository host; J private goals. Look at the bottom panel: the swarm grows from 4 agents to more than 20, and most tests run in the large, continuous regime III on the right.}
\label{fig:timeline}
\end{figure*}
